Adult attachment is commonly represented by attachment anxiety and avoidance, but these person-level dimensions do not by themselves specify how momentary relational conditions are perceived, compared with desired conditions, translated into action, and updated after interpersonal feedback. We present the Attachment Regulation and Control (ARC) framework as an integrative feedback-regulation account that makes these temporal assumptions explicit and empirically testable. In ARC, control refers to feedback-based regulation toward relational goals in the control-theoretic sense, not interpersonal dominance, coercion, or a generalized need for control. ARC does not claim a novel control-system architecture. Instead, it synthesizes established control-theoretic attachment ideas, secure-base representations, perceived partner responsiveness, relational discrepancy, dyadic regulation, and recovery into a parsimonious sequence for adult romantic relationships: perceived relational state, desired relational state, discrepancy, regulatory action, partner feedback, state updating, and recovery. The framework treats the discrepancy score itself as a falsifiable modeling restriction by comparing it with perceived-state-only, unconstrained perceived-plus-desired-state, and asymmetric discrepancy alternatives. Attachment anxiety and avoidance form the current person-level core and are modeled as moderators of channel-specific regulation. Positive secure-base representation is not yet required by the core; it is a prespecified candidate whose promotion depends on demonstrating incremental prediction of preregistered recovery outcomes beyond anxiety and avoidance. Mentalized affectivity and emotion-regulation flexibility remain optional moderators, and adult disorganization remains an independently measured extension outcome. The article derives core, candidate-core, and extension hypotheses, specifies competing observation models, and proposes staged intensive longitudinal and recovery-focused human tests.
